\documentclass[11pt]{article}
\usepackage[T1]{fontenc}
\usepackage{amsmath,amssymb,amsthm,lmodern}
\usepackage[margin=1in]{geometry}
\usepackage{graphicx,tikz-cd,hyperref}
\hypersetup{hidelinks,hypertexnames=false,pdftitle={EGNO: The Mac Lane strictness theorem, Theorem 2.8.5},pdfauthor={Pavel Etingof; Shlomo Gelaki; Dmitri Nikshych; Victor Ostrik}}
\newcommand{\cC}{\mathcal C}
\newcommand{\cD}{\mathcal D}
\newcommand{\one}{\mathbf 1}
\newcommand{\id}{\operatorname{id}}
\newcommand{\End}{\operatorname{End}}
\newcommand{\Hom}{\operatorname{Hom}}
\newcommand{\sourceloc}[1]{\par\smallskip\noindent\textit{Source: #1.}\par\nopagebreak[4]\smallskip}
\newcommand{\editorial}[1]{\par\smallskip\noindent\textsf{Editorial note.} #1\par\smallskip}
\newsavebox{\diagramBox}
\newcommand{\fitdiagram}[1]{\sbox{\diagramBox}{#1}\ifdim\wd\diagramBox>.92\linewidth\resizebox{.92\linewidth}{!}{\usebox{\diagramBox}}\else\usebox{\diagramBox}\fi}
\title{The Mac Lane Strictness Theorem\\Theorem 2.8.5}
\author{Pavel Etingof \and Shlomo Gelaki \and Dmitri Nikshych \and Victor Ostrik}
\date{Private reading transcription of the 2015 author final copy}
\begin{document}
\maketitle

% BEGIN CONTEXT EGNO:2.8.5
\noindent\textbf{Source and scope.}
\emph{Tensor Categories}, Mathematical Surveys and Monographs, volume 205,
American Mathematical Society, 2015, author final copy (362 physical pages),
\href{https://math.mit.edu/~etingof/egnobookfinal.pdf}{\nolinkurl{https://math.mit.edu/~etingof/egnobookfinal.pdf}}.
The primary selection is the original statement of Theorem 2.8.5 on physical
page 53 (printed page 36) and its complete printed proof on physical pages
54--55 (printed pages 37--38), ending before Remark 2.8.6.

The bibliographical notes on physical page 63 (printed page 46) say:
``The proof of the Mac Lane strictness Theorem 2.8.5 presented here is given
by Joyal and Street [JoyS5]. The coherence Theorem 2.9.2 is due to Mac Lane
[Mac1].'' Theorem 2.9.2 and its original short verbal proof are retained
below as uncounted support for this single primary target.

\editorial{The source's one-line essential-surjectivity observation,
strictness observation, and final reductions to pentagon and triangle
identities remain at their printed level of detail. The selected supporting
statements are background, not additional primary results. No proof step is
expanded here. The separately retained author corrections dated February 14,
2026 clarify the later Remark 2.8.6; that remark is outside this selection,
and its corrected tensor-product formula has not been substituted into the
proof of Theorem 2.8.5.}

\noindent\textbf{Original source footer.}
``Author's final version made available with permission of the publisher,
American Mathematical Society. License or copyright restrictions may apply
to redistribution; see \url{http://www.ams.org/publications/ebooks/terms}.''
This source transcription is kept private, in accordance with the author's
restriction on posting the manuscript on other websites.

\subsection*{Selected source conventions}
\sourceloc{Section 1.1, physical page 18, printed page 1}
The identity endofunctor of a category $\cC$ will be denoted $\id_{\cC}$.
We say that a functor $F:\cC\to\cD$ is an \emph{equivalence} if there is a
functor $F^{-1}:\cD\to\cC$, called a \emph{quasi-inverse} of $F$, such that
$F^{-1}\circ F\cong\id_{\cC}$ and $F\circ F^{-1}\cong\id_{\cD}$.

A category is called \emph{locally small} if for any objects $X,Y$,
$\Hom_{\cC}(X,Y)$ is a set, and is called \emph{essentially small} if in
addition its isomorphism classes of objects form a set. In other words,
an essentially small category is a category equivalent to a small category
(i.e., to one where objects and morphisms form a set). All categories
considered in this book will be locally small (except in the section on
2-categories), and most of them will be essentially small. So set-theoretical
subtleties will not play any role in this book.\footnote[1]{A well-known
foundational subtlety of category theory (which, luckily, does not play an
essential role in most category-theoretical considerations) is that in
general, objects and morphisms of a category are defined to form a class
rather than a set (this allows one to consider the category of all sets
without running into Russell's paradox---the non-existence of the set of
all sets). For example, the category of all sets or all vector spaces over
a field is not small or even essentially small (although it is locally
small). However, in this book we will usually consider categories with
finiteness conditions, such as the category of finite dimensional vector
spaces, so this issue is not going to be relevant most of the time. In fact,
even when we work with categories that are not essentially small (such as
the category of all vector spaces), we will allow ourselves to abuse
terminology and speak about ``the set of isomorphism classes of objects''
of such a category. This does not create a confusion in our context.}

\subsection*{Monoidal categories, associators, and unit constraints}
\sourceloc{Definition 2.1.1, physical pages 38--39, printed pages 21--22}
A monoidal category is a quintuple $(\cC,\otimes,a,\one,\iota)$ where
$\cC$ is a category, $\otimes:\cC\times\cC\to\cC$ is a bifunctor called
the tensor product bifunctor,
$a:(-\otimes-)\otimes-\xrightarrow{\sim}-\otimes(-\otimes-)$ is a natural
isomorphism:
\begin{equation}\tag{2.1}
a_{X,Y,Z}:(X\otimes Y)\otimes Z\xrightarrow{\sim}X\otimes(Y\otimes Z),
\qquad X,Y,Z\in\cC
\end{equation}
called the associativity constraint (or associativity isomorphism),
$\one\in\cC$ is an object of $\cC$, and
$\iota:\one\otimes\one\xrightarrow{\sim}\one$ is an isomorphism,
subject to the following two axioms.

\noindent\textbf{1. The pentagon axiom.} The diagram
\begin{equation}\tag{2.2}
\fitdiagram{\begin{tikzcd}[ampersand replacement=\&,column sep=large,row sep=large]
\& ((W\otimes X)\otimes Y)\otimes Z
 \arrow[ld,"{a_{W,X,Y}\otimes\id_Z}"']
 \arrow[rd,"{a_{W\otimes X,Y,Z}}"] \& \\
(W\otimes(X\otimes Y))\otimes Z
 \arrow[d,"{a_{W,X\otimes Y,Z}}"'] \& \&
(W\otimes X)\otimes(Y\otimes Z)
 \arrow[d,"{a_{W,X,Y\otimes Z}}"] \\
W\otimes((X\otimes Y)\otimes Z)
 \arrow[rr,"{\id_W\otimes a_{X,Y,Z}}"'] \& \& W\otimes(X\otimes(Y\otimes Z))
\end{tikzcd}}
\end{equation}
is commutative for all objects $W,X,Y,Z$ in $\cC$.

\noindent\textbf{2. The unit axiom.} The functors
\begin{align}\tag{2.3}
L_{\one}&:X\longmapsto\one\otimes X\quad\text{and}\\[-2pt]\tag{2.4}
R_{\one}&:X\longmapsto X\otimes\one
\end{align}
of left and right multiplication by $\one$ are autoequivalences of $\cC$.

\sourceloc{Section 2.2 opening and Definition 2.2.1, physical pages 39--40, printed pages 22--23}
Let $(\cC,\otimes,a,\one,\iota)$ be a monoidal category. Define natural isomorphisms
\begin{equation}\tag{2.5}
l_X:\one\otimes X\to X\qquad\text{and}\qquad r_X:X\otimes\one\to X
\end{equation}
in such a way that $L_{\one}(l_X)$ and $R_{\one}(r_X)$ are equal,
respectively, to the compositions
\begin{equation}\tag{2.6}
\one\otimes(\one\otimes X)\xrightarrow{a^{-1}_{\one,\one,X}}
(\one\otimes\one)\otimes X\xrightarrow{\iota\otimes\id_X}\one\otimes X,
\end{equation}
\begin{equation}\tag{2.7}
(X\otimes\one)\otimes\one\xrightarrow{a_{X,\one,\one}}
X\otimes(\one\otimes\one)\xrightarrow{\id_X\otimes\iota}X\otimes\one.
\end{equation}
Isomorphisms (2.5) are called the left and right unit constraints or unit isomorphisms.

\sourceloc{Proposition 2.2.3 and Proposition 2.2.4 (supporting statements), physical pages 40--41, printed pages 23--24}
The ``triangle'' diagram
\begin{equation}\tag{2.10}
\fitdiagram{\begin{tikzcd}[ampersand replacement=\&,column sep=large,row sep=large]
(X\otimes\one)\otimes Y \arrow[rr,"{a_{X,\one,Y}}"]
 \arrow[rd,"{r_X\otimes\id_Y}"'] \& \&
X\otimes(\one\otimes Y)\arrow[ld,"{\id_X\otimes l_Y}"]\\
\& X\otimes Y \&
\end{tikzcd}}
\end{equation}
is commutative for all $X,Y\in\cC$.

The following diagrams commute for all objects $X,Y\in\cC$:
\begin{equation}\tag{2.12}
\fitdiagram{\begin{tikzcd}[ampersand replacement=\&,column sep=large,row sep=large]
(\one\otimes X)\otimes Y\arrow[rr,"{a_{\one,X,Y}}"]
 \arrow[rd,"{l_X\otimes\id_Y}"'] \& \&
\one\otimes(X\otimes Y)\arrow[ld,"{l_{X\otimes Y}}"]\\
\& X\otimes Y \&
\end{tikzcd}}
\end{equation}
\begin{equation}\tag{2.13}
\fitdiagram{\begin{tikzcd}[ampersand replacement=\&,column sep=large,row sep=large]
(X\otimes Y)\otimes\one\arrow[rr,"{a_{X,Y,\one}}"]
 \arrow[rd,"{r_{X\otimes Y}}"'] \& \&
X\otimes(Y\otimes\one)\arrow[ld,"{\id_X\otimes r_Y}"]\\
\& X\otimes Y \&
\end{tikzcd}}
\end{equation}

\subsection*{Monoidal functors and their unit compatibility}
\sourceloc{Definition 2.4.1, physical page 47, printed page 30}
Let $(\cC,\otimes,\one,a,\iota)$ and
$(\cC',\otimes',\one',a',\iota')$ be two monoidal categories.
A monoidal functor from $\cC$ to $\cC'$ is a pair $(F,J)$, where
$F:\cC\to\cC'$ is a functor, and
\begin{equation}\tag{2.22}
J_{X,Y}:F(X)\otimes'F(Y)\xrightarrow{\sim}F(X\otimes Y)
\end{equation}
is a natural isomorphism, such that $F(\one)$ is isomorphic to $\one'$ and the diagram
\begin{equation}\tag{2.23}
\fitdiagram{\begin{tikzcd}[ampersand replacement=\&,column sep=huge,row sep=large]
(F(X)\otimes'F(Y))\otimes'F(Z)
 \arrow[r,"{a'_{F(X),F(Y),F(Z)}}"]
 \arrow[d,"{J_{X,Y}\otimes'\id_{F(Z)}}"'] \&
F(X)\otimes'(F(Y)\otimes'F(Z))
 \arrow[d,"{\id_{F(X)}\otimes'J_{Y,Z}}"]\\
F(X\otimes Y)\otimes'F(Z)
 \arrow[d,"{J_{X\otimes Y,Z}}"'] \&
F(X)\otimes'F(Y\otimes Z)
 \arrow[d,"{J_{X,Y\otimes Z}}"]\\
F((X\otimes Y)\otimes Z)
 \arrow[r,"{F(a_{X,Y,Z})}"'] \& F(X\otimes(Y\otimes Z))
\end{tikzcd}}
\end{equation}
is commutative for all $X,Y,Z\in\cC$ (``the monoidal structure axiom'').

A monoidal functor $F$ is said to be an \emph{equivalence of monoidal categories}
if it is an equivalence of ordinary categories.

\sourceloc{Canonical unit isomorphism and Proposition 2.4.3, physical pages 47--48, printed pages 30--31}
It turns out that if $F$ is a monoidal functor, then there is a canonical
isomorphism $\varphi:\one'\to F(\one)$. This isomorphism is defined by the
commutative diagram
\begin{equation}\tag{2.24}
\fitdiagram{\begin{tikzcd}[ampersand replacement=\&,column sep=huge,row sep=large]
\one'\otimes'F(\one)\arrow[r,"{l'_{F(\one)}}"]
 \arrow[d,"{\varphi\otimes'\id_{F(\one)}}"'] \& F(\one)\arrow[d,"{F(l_{\one})^{-1}}"]\\
F(\one)\otimes'F(\one)\arrow[r,"{J_{\one,\one}}"'] \& F(\one\otimes\one)
\end{tikzcd}}
\end{equation}
where $l,r,l',r'$ are the unit isomorphisms for $\cC$ and $\cC'$ defined in (2.5).

For any monoidal functor $(F,J):\cC\to\cC'$, the diagrams
\begin{equation}\tag{2.25}
\fitdiagram{\begin{tikzcd}[ampersand replacement=\&,column sep=huge,row sep=large]
\one'\otimes'F(X)\arrow[r,"{l'_{F(X)}}"]
 \arrow[d,"{\varphi\otimes'\id_{F(X)}}"'] \& F(X)\arrow[d,"{F(l_X)^{-1}}"]\\
F(\one)\otimes'F(X)\arrow[r,"{J_{\one,X}}"'] \& F(\one\otimes X)
\end{tikzcd}}
\end{equation}
and
\begin{equation}\tag{2.26}
\fitdiagram{\begin{tikzcd}[ampersand replacement=\&,column sep=huge,row sep=large]
F(X)\otimes'\one'\arrow[r,"{r'_{F(X)}}"]
 \arrow[d,"{\id_{F(X)}\otimes'\varphi}"'] \& F(X)\arrow[d,"{F(r_X)^{-1}}"]\\
F(X)\otimes'F(\one)\arrow[r,"{J_{X,\one}}"'] \& F(X\otimes\one)
\end{tikzcd}}
\end{equation}
are commutative for all $X\in\cC$.

\subsection*{Endofunctors and strictness}
\sourceloc{Example 2.3.12, physical page 45, printed page 28; selected sentences}
Let $\cC$ be a category. Then the category $\End(\cC)$ of all functors from
$\cC$ to itself is a monoidal category, where $\otimes$ is given by composition
of functors. The associativity isomorphism in this category is the identity.
The unit object is the identity functor, and the structure morphisms are obvious.

\sourceloc{Definition 2.8.1 and Example 2.8.2, physical page 53, printed page 36}
A monoidal category $\cC$ is \emph{strict} if for all objects $X,Y,Z$ in $\cC$
one has equalities $(X\otimes Y)\otimes Z=X\otimes(Y\otimes Z)$ and
$X\otimes\one=X=\one\otimes X$, and the associativity and unit constraints
are the identity maps.
The category $\End(\cC)$ of endofunctors of a category $\cC$ (see Example 2.3.12) is strict.
% END CONTEXT EGNO:2.8.5

\clearpage
\sourceloc{Theorem 2.8.5, physical page 53, printed page 36}
% BEGIN PRIMARY STATEMENT EGNO:2.8.5
\noindent\textbf{Theorem 2.8.5.}
\emph{Any monoidal category is monoidally equivalent to a strict monoidal category.}
% END PRIMARY STATEMENT EGNO:2.8.5

\sourceloc{Complete printed proof, physical pages 54--55, printed pages 37--38}
% BEGIN PRIMARY PROOF EGNO:2.8.5
\begin{proof}
We will establish an equivalence between $\cC$ and the monoidal category
of right $\cC$-module endofunctors of $\cC$, which we will discuss in more
detail in Chapter 7. The non-categorical counterpart of this result is the fact
that every monoid $M$ is isomorphic to the monoid consisting of maps from
$M$ to itself commuting with the right multiplication.

For a monoidal category $\cC$ let $\cC'$ be the monoidal category defined
as follows. The objects of $\cC'$ are pairs $(F,c)$ where $F:\cC\to\cC$
is a functor and
\[
c_{X,Y}:F(X)\otimes Y\xrightarrow{\sim}F(X\otimes Y)
\]
is a natural isomorphism such that the following diagram is commutative for
all objects $X,Y,Z$ in $\cC$:
\begin{equation}\tag{2.38}
\fitdiagram{\begin{tikzcd}[ampersand replacement=\&,column sep=large,row sep=large]
\& (F(X)\otimes Y)\otimes Z
 \arrow[ld,"{c_{X,Y}\otimes\id_Z}"']
 \arrow[rd,"{a_{F(X),Y,Z}}"] \& \\
F(X\otimes Y)\otimes Z\arrow[d,"{c_{X\otimes Y,Z}}"'] \& \&
F(X)\otimes(Y\otimes Z)\arrow[d,"{c_{X,Y\otimes Z}}"]\\
F((X\otimes Y)\otimes Z)\arrow[rr,"{F(a_{X,Y,Z})}"'] \& \&
F(X\otimes(Y\otimes Z)).
\end{tikzcd}}
\end{equation}
A morphism $\theta:(F^1,c^1)\to(F^2,c^2)$ in $\cC'$ is a natural
transformation $\theta:F^1\to F^2$ such that the following square commutes
for all objects $X,Y$ in $\cC$:
\begin{equation}\tag{2.39}
\fitdiagram{\begin{tikzcd}[ampersand replacement=\&,column sep=huge,row sep=large]
F^1(X)\otimes Y\arrow[r,"{c^1_{X,Y}}"]
 \arrow[d,"{\theta_X\otimes\id_Y}"'] \&F^1(X\otimes Y)\arrow[d,"{\theta_{X\otimes Y}}"]\\
F^2(X)\otimes Y\arrow[r,"{c^2_{X,Y}}"'] \&F^2(X\otimes Y),
\end{tikzcd}}
\end{equation}
Composition of morphisms is the vertical composition of natural
transformations. The tensor product of objects is given by
$(F^1,c^1)\otimes(F^2,c^2)=(F^1F^2,c)$ where $c$ is given by a composition
\begin{equation}\tag{2.40}
\fitdiagram{\begin{tikzcd}[ampersand replacement=\&,column sep=large]
F^1F^2(X)\otimes Y\arrow[r,"{c^1_{F^2(X),Y}}"] \&
F^1(F^2(X)\otimes Y)\arrow[r,"{F^1(c^2_{X,Y})}"] \& F^1F^2(X\otimes Y)
\end{tikzcd}}
\end{equation}
for all $X,Y\in\cC$, and the tensor product of morphisms is the composition
of natural transformations. Thus $\cC'$ is a strict tensor category (the unit
object is the identity functor).

Consider a functor of left multiplication $L:\cC\to\cC'$ given by
\[
L(X)=(X\otimes-,a_{X,-,-}),\qquad L(f)=(f\otimes-).
\]
Note that the diagram (2.38) for $L$ is nothing but the pentagon diagram (2.2).

We will show that this functor $L$ is a monoidal equivalence. First of all,
note that any $(F,c)$ in $\cC'$ is isomorphic to $L(F(\one))$.

Let us now show that $L$ is fully faithful. Let $\theta:L(X)\to L(Y)$ be a
morphism in $\cC'$. Define $f:X\to Y$ to be the composition
\begin{equation}\tag{2.41}
X\xrightarrow{r_X^{-1}}X\otimes\one\xrightarrow{\theta_{\one}}
Y\otimes\one\xrightarrow{r_Y}Y,
\end{equation}
where $r$ is the unit constraint. We claim that for all $Z$ in $\cC$ one has
$\theta_Z=f\otimes\id_Z$ (so that $\theta=L(f)$ and $L$ is full).
Indeed, this follows from the commutativity of the following diagram
\begin{equation}\tag{2.42}
\fitdiagram{\begin{tikzcd}[ampersand replacement=\&,column sep=large,row sep=large]
X\otimes Z\arrow[r,"{r_X^{-1}\otimes\id_Z}"]
 \arrow[d,"{f\otimes\id_Z}"'] \&
(X\otimes\one)\otimes Z\arrow[r,"{a_{X,\one,Z}}"]
 \arrow[d,"{\theta_{\one}\otimes\id_Z}"] \&
X\otimes(\one\otimes Z)\arrow[r,"{X\otimes l_Z}"]
 \arrow[d,"{\theta_{\one\otimes Z}}"] \&X\otimes Z\arrow[d,"{\theta_Z}"]\\
Y\otimes Z\arrow[r,"{r_Y^{-1}\otimes\id_Z}"'] \&
(Y\otimes\one)\otimes Z\arrow[r,"{a_{Y,\one,Z}}"'] \&
Y\otimes(\one\otimes Z)\arrow[r,"{Y\otimes l_Z}"'] \&Y\otimes Z,
\end{tikzcd}}
\end{equation}
where the rows are the identity morphisms by the triangle axiom (2.10),
the left square commutes by the definition of $f$, the right square commutes
by the naturality of $\theta$, and the central square commutes since $\theta$
is a morphism in $\cC'$.

Next, if $L(f)=L(g)$ for some morphisms $f,g$ in $\cC$ then
$f\otimes\id_{\one}=g\otimes\id_{\one}$ so that $f=g$ by the naturality of $r$.
So $L$ is faithful. Thus, it is an equivalence.

Finally, we define a monoidal functor structure
\[
\phi:\one_{\cC'}\xrightarrow{\sim}L(\one_{\cC})
\quad\text{and}\quad
J_{X,Y}:L(X)\circ L(Y)\xrightarrow{\sim}L(X\otimes Y)
\]
on $L$ by $\phi=l^{-1}:(\id_{\cC},\id)\xrightarrow{\sim}
(\one\otimes-,a_{\one,-,-})$ and
\begin{gather*}
J_{X,Y}=a^{-1}_{X,Y,-}:\bigl((X\otimes(Y\otimes-)),
(\id_X\otimes a_{Y,-,-}\circ a_{X,Y\otimes-,-})\bigr)\\
\xrightarrow{\sim}\bigl((X\otimes Y)\otimes-,a_{X\otimes Y,-,-}\bigr).
\end{gather*}
The diagram (2.39) for the latter natural isomorphism is the pentagon
diagram in $\cC$. For the functor $L$ the hexagon diagram (2.23) in the
definition of a monoidal functor also reduces to the pentagon diagram in
$\cC$. The square diagrams (2.25) and (2.26) reduce to triangles, one of which
is the triangle axiom (2.10) for $\cC$ and another is (2.13).
\end{proof}
% END PRIMARY PROOF EGNO:2.8.5

\subsection*{Uncounted support: the coherence theorem}
\sourceloc{Theorem 2.9.2 and its complete printed verbal proof, physical page 57, printed page 40}
% BEGIN SUPPORT STATEMENT EGNO:2.9.2
\noindent\textbf{Theorem 2.9.2. (Coherence Theorem)}
\emph{Let $X_1,\ldots,X_n\in\cC$. Let $P_1,P_2$ be any two parenthesized
products of $X_1,\ldots,X_n$ (in this order) with arbitrary insertions of the
unit object $\one$. Let $f,g:P_1\to P_2$ be two isomorphisms, obtained by
composing associativity and unit isomorphisms and their inverses possibly
tensored with identity morphisms. Then $f=g$.}
% END SUPPORT STATEMENT EGNO:2.9.2

% BEGIN SUPPORT PROOF EGNO:2.9.2
\begin{proof}
We derive this theorem as a corollary of the Mac Lane's strictness
Theorem 2.8.5. Let $L:\cC\to\cC'$ be a monoidal equivalence between $\cC$
and a strict monoidal category $\cC'$. Consider a diagram in $\cC$
representing $f$ and $g$ and apply $L$ to it. Over each arrow of the resulting
diagram representing an associativity isomorphism, let us build a rectangle
as in (2.23), and do similarly for the unit morphisms. This way we obtain a
prism one of whose faces consists of identity maps (associativity and unit
isomorphisms in $\cC'$) and whose sides are commutative. Hence, the other
face is commutative as well, i.e., $f=g$.
\end{proof}
% END SUPPORT PROOF EGNO:2.9.2
\end{document}
